% ============================================
% Supplementary Note 4: Nutrient Dependence of Pairwise Invasion Resistance
% ============================================
\subsection*{Supplementary Note 4: Nutrient dependence of pairwise invasion resistance}
\label{sec:suppl:note4}

To operationally estimate pairwise invasion resistance, we performed reciprocal invasion assays using the 12 most abundant isolates (ranked by their corresponding ASV abundances across the sequencing dataset), testing each pair in both directions to assess competitive outcomes and detect potential bistability. Full assay protocol (95:5 resident:invader ratio, 7-cycle propagation, colony counting, and classification thresholds) is described in Supplementary Methods.

Invasion assays were performed across all three nutrient conditions: Nutr$-$, Base, and Nutr$+$. We use failed invasion as the operational assay readout for competitive exclusion, defined as cases in which the introduced species remained below 1\% relative abundance after propagation. Among the assayed abundant isolates, the frequency of failed invasion increased with nutrient concentration, consistent with the coalescence outcome patterns observed in community experiments (Supplementary Figs.~7--9). Under Nutr$-$, more pairs achieved stable coexistence, while Nutr$+$ conditions led to more frequent exclusion of one species by the other. Higher failed-invasion frequency among these assayed isolates is consistent with stronger resistance to cross-community establishment~\citep{Hu2025,Kurkjian2021} and with prior work linking nutrient and interaction strength to microbial competitive dynamics~\citep{Hu2022,Ratzke2020}, which can generate more correlated parental-community fates after coalescence. This provides a plausible link by which pairwise invasion resistance could favor Dominance rather than balanced Mixture, although these assays alone do not establish causality beyond the evidence shown in the figures.
